Information sharing changes both posterior quality and the diversity of actions
available for discovery. We study that joint effect in finite one-hit search.
First, exact channel comparisons show that one-person Bayes accuracy does not
determine the value of a pooled action budget. Second, an incremental-sharing
identity separates the accuracy gained by enlarging a pooled block from the
independent rescue action that pooling removes. This yields an exact local
criterion: sharing helps precisely when pooled residual error contracts faster
than the removed private-failure factor. A centralized posterior top-$L$
portfolio can recover the declared private baseline at full action capacity,
but this is an authority benchmark, not a decentralized implementation.

We then study a two-agent equal-split discovery game with a hidden mixture of
common and conditionally independent signal sources. At signal accuracy
$p=3/5$, sharing strictly improves discovery and each symmetric agent's payoff
over the private anonymous symmetric Bayes--Nash selection exactly when
$\rho\in((5\sqrt{73}-17)/48,1)$. This theorem is selection-dependent: the
shared outcome is the declared anonymous posterior-only, provenance-blind,
identical-mixing equilibrium, not an every-equilibrium result. Opposite
constant-target private equilibria exist and attain discovery one, while
ownership-aware symmetric strategies can split targets on public disagreement.
Sharing reveals the agreement pattern and changes posterior beliefs about
dependence; it does not reveal the realized common-versus-independent source
branch. The combined evidence is analytic and exact finite computation. It
uses no human or real data, and the bounded channel registry is not a universal
information order.
